\begin{statementbox}

	\textbf{\textcolor{CStatement}{条件}}
	\begin{description}[style=nextline, leftmargin=2.8em, labelsep=0.5em, itemsep=0.35em]
		\item[\textbf{\textcolor{CStatement}{条件1（十字分解）：}}] 对于二次三项式 $x^2+px+q$，存在实数 $a,b$ 使得 $p=a+b$ 且 $q=ab$。
	\end{description}

	\textbf{\textcolor{CStatement}{结论}}
	\begin{description}[style=nextline, leftmargin=2.8em, labelsep=0.5em, itemsep=0.45em]
		\item[\textbf{\textcolor{CStatement}{结论一（十字分解式）：}}]
			\textbf{\textcolor{CStatement}{直观描述：}}
			将满足条件的 $a,b$ 直接写入因式，得到乘积形式。
			\[
				x^2+px+q=(x+a)(x+b).
			\]

		\item[\textbf{\textcolor{CStatement}{常见变形（根式分解）：}}]
			\textbf{\textcolor{CStatement}{直观描述：}}
			若二次方程有实根，可用根直接写出分解。
			\[
				x^2+px+q=(x-x_1)(x-x_2),
			\]
			其中 $x_1,x_2$ 是方程 $x^2+px+q=0$ 的两根。
	\end{description}

	\medskip
	\textbf{\textcolor{CStatement}{使用提醒：}}
	\begin{itemize}[leftmargin=*, itemsep=0.2em, topsep=0.2em]
		\item 该方法特别适用于 $p,q$ 为整数或简单有理数的情形，试根时可先枚举 $q$ 的因子对。
		\item 若 $a,b$ 为整数，则分解式可直接看出；若为无理数，则一般用求根公式替代。
	\end{itemize}
\end{statementbox}
